\documentclass[aspectratio=169,12pt]{beamer}
\usepackage[T1]{fontenc}
\usepackage[utf8]{inputenc}
\usepackage[brazilian]{babel}
\usepackage{helvet}
\urlstyle{same}
\usepackage{booktabs,tikz,array,xcolor,pgfplots}
\usetikzlibrary{arrows.meta,positioning}
\pgfplotsset{compat=1.18}
\input{../.build/tcc/numbers.tex}
\input{../.build/tcc/scaling-coordinates.tex}
\newcommand{\numpt}[1]{\pgfmathprintnumber[fixed,precision=2,use comma,1000 sep={\,}]{#1}}
\newcommand{\latpt}[1]{\pgfmathprintnumber[fixed,precision=3,use comma,1000 sep={\,}]{#1}}
\newcommand{\ratiopt}[1]{\pgfmathprintnumber[fixed,precision=1,use comma,1000 sep={\,}]{#1}}
\pgfplotsset{every axis/.append style={
 tick label style={font=\scriptsize,/pgf/number format/use comma},
 label style={font=\small},legend style={font=\scriptsize}}}
\newcolumntype{L}[1]{>{\raggedright\arraybackslash}p{#1}}
\definecolor{ink}{HTML}{142B42}
\definecolor{accent}{HTML}{087F72}
\definecolor{muted}{HTML}{566575}
\setbeamercolor{normal text}{fg=ink,bg=white}
\setbeamercolor{structure}{fg=accent}
\setbeamercolor{frametitle}{fg=ink}
\setbeamerfont{frametitle}{size=\LARGE,series=\bfseries}
\setbeamerfont{title}{size=\Huge,series=\bfseries}
\setbeamertemplate{frametitle}[default][left,leftskip=0.7cm,rightskip=0.7cm]
\setbeamertemplate{navigation symbols}{}
\setbeamertemplate{footline}{\hfill\color{muted}\small\insertframenumber\hspace{10mm}\vspace{4mm}}
\setbeamertemplate{itemize item}{\color{accent}\small\textbullet}
\setbeamersize{text margin left=10mm,text margin right=10mm}
\newcommand{\source}[1]{\par\vspace{4mm}{\fontsize{9}{11}\selectfont\color{muted}#1}}
\newcommand{\gap}{\par\vspace{5mm}}
\tikzset{stage/.style={draw=accent,rounded corners=2pt,align=center,minimum height=10mm,text width=30mm},>=Stealth}
\title{JayDB2}
\subtitle{Transações acompanhadas de provas\\de conformidade com as leis do administrador}
\author{Gabriel Jota Lizardo}
\institute{Pontifícia Universidade Católica de Minas Gerais}
\date{2026}
\hypersetup{pdftitle={JayDB2: transações acompanhadas de provas},
 pdfauthor={Gabriel Jota Lizardo},pdflang={pt-BR}}
\begin{document}

% 1
\begin{frame}[plain]
\vspace{5mm}
{\Huge\bfseries JayDB2}\par\gap
{\LARGE Transações acompanhadas de provas}\par\vspace{3mm}
{\large Conformidade com as leis do administrador}\par\vspace{10mm}
Gabriel Jota Lizardo\par
{\large PUC Minas\hspace{6mm}2026}
\note{O JayDB2 é um banco de dados experimental escrito em Bend 2, com um host em Rust. O cliente envia um programa junto com uma prova de que ele respeita as leis instaladas pelo administrador. A pergunta central é como conectar essa prova à execução e à publicação da transação. Vou apresentar a interface, os teoremas sobre o código executável, duas aplicações certificadas e o protocolo de avaliação do custo dessa certificação. A base de confiança fica explícita ao longo da construção.}
\end{frame}

% 2
\begin{frame}{Permissões e integridade}
\centering
\begin{tikzpicture}[node distance=7mm,font=\small]
\node[stage] (a) {Alice disponível\\índice: ready};
\node[stage,right=of a] (b) {Manutenção\\índice transitório};
\node[stage,right=of b] (c) {Alice indisponível\\índice: busy};
\draw[->,thick] (a) -- (b);
\draw[->,thick] (b) -- (c);
\draw[accent,thick] ([yshift=-6mm]a.south west) -- ([yshift=-6mm]c.south east);
\node[below=9mm of b,align=center] {Cada acesso pertence a Alice};
\end{tikzpicture}
\gap
O índice pode ficar incompleto na proposta privada.\\
A publicação exige a projeção exata da linha.
\source{Aplicação: applications/indexed}
\note{Alice e Bob têm uma linha de disponibilidade e um índice próprios. Para mudar a disponibilidade, o programa remove a entrada antiga, grava a nova e altera a linha principal. Durante esses passos o índice pode ficar incompleto, mas cada acesso continua restrito às tabelas de Alice. Ao publicar sucesso, o índice precisa ser exatamente a projeção da linha. Ler o índice de Bob já viola a permissão, mesmo que a consulta descarte o resultado ou recuse depois. Esse exemplo separa uma obrigação sobre cada prefixo de outra sobre a publicação.}
\end{frame}

% 3
\begin{frame}{Perguntas de pesquisa}
\begin{tabular}{@{}L{.13\linewidth}L{.79\linewidth}@{}}
\textbf{RQ1} & Quais garantias decorrem do código executável, sob qual base de confiança?\\[6mm]
\textbf{RQ2} & Certificados combinam acesso em cada prefixo e invariantes de publicação?\\[6mm]
\textbf{RQ3} & Quanto custa certificar e executar no serviço atual?\\
\end{tabular}
\note{A primeira pergunta exige teoremas sobre as funções que realmente decidem admissão, execução e publicação. A segunda trata da expressão de leis de aplicação, com controles inválidos que falham na obrigação correta. A terceira mede o custo do caminho ordinário e o que muda quando um programa preparado reutiliza a certificação. O método combina provas, controles negativos, execução nativa e experimentos com fronteiras de medição definidas. As medições quantitativas desta versão ainda estão pendentes e aparecem em vermelho.}
\end{frame}

% 4
\begin{frame}{Interface de programas e certificados}
\begin{tabular}{@{}L{.25\linewidth}L{.68\linewidth}@{}}
\textbf{Administrador} & Esquema e leis sobre situações.\\[5mm]
\textbf{Cliente} & Credencial, programa puro, argumentos e prova.\\[5mm]
\textbf{JayDB2} & Autenticação atual, admissão exata e execução por ações.\\
\end{tabular}
\gap
{\color{accent}Todas as leis instaladas participam do certificado.}
\source{Situações: estado estável, prefixo e encerramento}
\note{O administrador define as proposições da aplicação. O programa recebe o contexto autenticado e os argumentos, e produz ações restritas. Somente o interpretador fornece resultados de operações sobre os dados. O cliente não recebe o estado inteiro. O certificado quantifica sobre todos os estados válidos compatíveis com o contexto e sobre todos os prefixos de execução. Assim, ele não depende de revelar o estado privado ao compilador. A identidade vem da credencial verificada no registro atual. A terminação é uma obrigação adicional da admissão. Fontes: contract/Semantics.bend e contract/Requirements.bend.}
\end{frame}

% 5
\begin{frame}{Execução e publicação}
\centering
\begin{tikzpicture}[font=\small,node distance=8mm and 13mm]
\node[stage] (root) {Estado publicado\\$s$};
\node[stage,right=of root] (work) {Execução privada\\$w$};
\node[stage,right=of work] (finish) {Finalização\\resposta exata};
\node[below=12mm of finish,align=center] (commit) {\textbf{Commit}\\publicar $w$};
\node[below=12mm of work,align=center] (keep) {\textbf{Keep}\\preservar $s$};
\draw[->,thick] (root) -- (work);
\draw[->,thick] (work) -- (finish);
\draw[->,accent,thick] (finish) -- node[right]{sucesso} (commit);
\draw[->,thick] ([yshift=-7mm]finish.south) -| node[pos=.25,above,font=\scriptsize]{demais resultados} (keep.north);
\end{tikzpicture}
\gap
Rust obedece à decisão tipada produzida em Bend.
\note{A execução começa com uma proposta privada igual ao estado publicado. Leituras posteriores veem alterações anteriores da mesma chamada. O finalizador verifica estrutura, tamanho de estado e resposta e produz os bytes completos. A decisão de Commit ou Keep também é uma função Bend verificada. Sucesso publica a raiz selecionada, que pode ser igual à inicial em uma leitura. Os demais resultados preservam a raiz anterior. O host recolhe a saída completa e termina a limpeza antes de publicar e responder. Depois do commit, perder a resposta não desfaz a publicação. Fontes: src/Gate.bend, src/Execution.bend e proof/WorkerDecision.bend.}
\end{frame}

% 6
\begin{frame}{Arquitetura do caminho ordinário}
\centering
\begin{tikzpicture}[font=\small,node distance=9mm and 10mm,
 stage/.append style={text width=34mm}]
\node[stage] (client) {Cliente\\programa e prova};
\node[stage,right=of client] (host) {Rust\\quadro e raiz atual};
\node[stage,right=of host] (check) {Admissão\\dois modos e C};
\node[stage,below=of check] (worker) {Worker Bend\\Gate\\interpretador};
\node[stage,left=of worker] (decision) {Decisão verificada\\Commit / Keep};
\node[stage,left=of decision] (pub) {Rust\\publicação e resposta};
\draw[->,thick] (client) -- (host);
\draw[->,thick] (host) -- (check);
\draw[->,thick] (check) -- (worker);
\draw[->,thick] (worker) -- (decision);
\draw[->,thick] (decision) -- (pub);
\end{tikzpicture}
\source{Requisições FIFO e seriais. Host Rust: \numpt{\ServerNonblankLines} linhas não vazias.}
\note{Rust coordena processos, pipes e a raiz publicada. Node prepara fontes e chama as ferramentas oficiais. O Book verificado é o mesmo enviado ao compilador Bend para C, e Clang compila esse C sem reescrita. O worker chama Gate.begin, avança o interpretador e termina por Gate.finish ou Gate.early. WorkerBridge relaciona essa execução à composição dos teoremas. Cada Book gerado também prova generated\_oracle\_sound para o seu próprio oráculo de admissão. Assim, o oráculo aceito constrói exatamente o programa certificado. As requisições são processadas em série. O paralelismo interno do runtime não cria transações concorrentes.}
\end{frame}

% 7
\begin{frame}{Garantias do Core}
\begin{tabular}{@{}L{.27\linewidth}L{.66\linewidth}@{}}
\textbf{Admitted} & Identidade atual e início do programa exato.\\[4mm]
\textbf{Safe} & Todas as leis em cada prefixo real.\\[4mm]
\textbf{Completed} & Finalização, raiz válida e bytes exatos.\\[4mm]
\textbf{Progressing} & Programa aceito alcança um resultado final.\\
\end{tabular}
\source{Premissas comuns: configuração válida, estado válido e solidez do oráculo.}
\note{As quatro cláusulas falam da mesma entrada computada por Gate.begin. Admitted fixa a identidade e a máquina iniciada, Safe cobre todos os seus prefixos, Completed identifica o finalizador e a publicação exatos, e Progressing fornece um prefixo final finito. Se a proposta de sucesso ou recusa é preservada pela finalização, a resposta mantém esse resultado. Erros de dados continuam possíveis. Ingress estende o resultado à decodificação, UnifiedCore à seleção direta ou registrada, e os folds preservam validade por sequências de transações e controles. A lista completa está no apêndice. Fonte: contract/Theorems.bend.}
\end{frame}

% 8
\begin{frame}{Uma lição sobre contratos vazios}
\begin{columns}[T,onlytextwidth]
\column{.46\textwidth}
\textbf{Contrato anterior}\par\vspace{4mm}
Relação abstrata de ambiente.\par\vspace{3mm}
Um ambiente sem execuções úteis podia satisfazer as implicações.
\column{.48\textwidth}
\textbf{Contrato executável}\par\vspace{4mm}
Entrada e resultado computados.\par\vspace{3mm}
Aceitação exige progresso real.
\end{columns}
\gap
{\color{accent}Controle positivo: autenticação, admissão e sucesso com mudança de estado.}
\note{O contrato anterior deixava uma relação abstrata decidir quais execuções existiam. Um ambiente trivial podia tornar suas implicações verdadeiras sem executar uma transação útil. A correção foi formular os teoremas diretamente sobre Gate, o interpretador e o finalizador. O fixture positivo de tarefas usa política, programa, prova e semente reais e demonstra sucesso com mudança da raiz. Refutações excluem identidade falsa, início sem certificado, publicação em recusa e rejeição de um programa aceito pelo oráculo. Centenas de controles negativos precisam falhar exatamente na definição false, nos dois modos oficiais. Fontes: tests/proof/Positive.bend e tests/proof/Refutations.bend.}
\end{frame}

% 9
\begin{frame}{Base de confiança e alcance}
\begin{tabular}{@{}L{.12\linewidth}L{.81\linewidth}@{}}
\textbf{T1} & Correção dos verificadores, vínculo fonte/Book/artefato e pureza.\\[4mm]
\textbf{T2} & Preservação pelo compilador Bend--C e runtime.\\[4mm]
\textbf{T3} & Enquadramento, processos, bytes atuais e publicação física.\\
\end{tabular}
\gap
\small Fora dos teoremas: durabilidade e recuperação de falhas, falhas do SO,
recursos físicos e completude do verificador.
\source{Deveres T3 numerados em proof/README.md e detalhados no apêndice.}
\note{T1 inclui a soundness dos verificadores, a auditoria de pureza e a ligação do programa exato ao Book compilado. A inicialização depende da declaração exata de validade da semente admitida antes de executar. A prova do oráculo gerado já é uma proposição verificada. T2 preserva as transições ao compilar e executar. T3 especifica o que o host físico precisa fazer, inclusive encaminhar o registro preparado mais recente sem alterar bytes, obedecer Commit e Keep e classificar erros de ferramentas. Os teoremas cobrem o serviço lógico serial. Durabilidade, falhas do sistema operacional, recursos físicos e aceitação de toda prova matemática ficam fora desse alcance.}
\end{frame}

% 10
\begin{frame}{Programas preparados}
\centering
\begin{tikzpicture}[font=\small,node distance=9mm and 9mm]
\node[stage] (proof) {Prova universal\\$\forall c,a$};
\node[stage,right=of proof] (prepare) {Verificar e compilar\\uma vez};
\node[stage,right=of prepare] (handle) {Handle\\da instância};
\draw[->,thick] (proof) -- (prepare);
\draw[->,thick] (prepare) -- (handle);
\node[below=15mm of prepare,align=center,text width=105mm] (call) {Cada chamada: autenticação atual e argumentos atuais\\seleção verificada e execução privada};
\draw[->,thick] (handle.south) -- ++(0,-6mm) -| (call.north);
\end{tikzpicture}
\gap
Preparar e invalidar usam o controlador Bend.\\
Um handle revogado permanece ausente na instância.
\note{A preparação verifica um certificado universal e compila o programa uma única vez. Cada chamada por handle autentica no estado atual e especializa essa prova para o contexto e os argumentos recebidos. Não há verificador ou compilador na reutilização. O controlador verifica a admissão universal exata, o formato e a novidade vitalícia do handle e o tamanho do recibo antes de mudar o registro. Control e ServiceFold cobrem essas transições, e ServiceNoRevival mantém um handle removido ausente. Rust encaminha os últimos bytes do registro sem alteração. Handles desaparecem no reinício. A preparação concede uma referência ao programa, mantendo a identidade atual de cada usuário.}
\end{frame}

% 11
\begin{frame}{Aplicação: propriedade e índice exato}
\begin{tabular}{@{}L{.34\linewidth}L{.26\linewidth}L{.31\linewidth}@{}}
\toprule
Programa & Acesso & Índice ao publicar\\\midrule
Manutenção correta & Permitido & Exato\\[3mm]
Manutenção omitida & Permitido & Prova rejeitada\\[3mm]
Valor de índice falso & Permitido & Prova rejeitada\\[3mm]
Índice de Bob por Alice & Prova rejeitada & Sem execução\\
\bottomrule
\end{tabular}
\gap
Dois controles preservam a prova de permissão e falham na igualdade final.
\source{Policy.Rule e Policy.IndexRule são leis independentes.}
\note{A aplicação tem dois usuários fixos e uma linha booleana por usuário. O índice é uma tabela comum. A prova da ação segue os resultados reais de Data.operate e termina com a igualdade entre índice e projeção. As versões que omitem a manutenção ou gravam um valor falso têm provas independentes de permissão válidas para todos os estados e prefixos. O certificado completo falha exatamente na igualdade final. A leitura do índice de Bob por Alice falha na primeira obrigação de permissão. O host rejeita os três casos e consultas posteriores confirmam o estado preservado. Isso responde à RQ2 distinguindo duas causas de erro.}
\end{frame}

% 12
\begin{frame}{Aplicação: vouchers e recibos}
\begin{tabular}{@{}L{.34\linewidth}L{.59\linewidth}@{}}
\textbf{Primeira emissão} & Grava voucher e recibo na mesma transação.\\[4mm]
\textbf{Repetição idêntica} & Revalida permissão e retorna o resultado retido.\\[4mm]
\textbf{Entrada conflitante} & Recusa.\\[4mm]
\textbf{Cancelamento} & Faz as escritas privadas e recusa, sem publicação.\\
\end{tabular}
\source{A retirada explícita remove voucher e recibo e encerra a proteção do ID.}
\note{Se o serviço publica e a resposta se perde, repetir uma emissão sem proteção pode duplicar seu efeito. Aqui o ID do pedido fica ligado à operação, à entrada exata e ao resultado. A consulta ao recibo depende da identidade e da permissão atuais. As leis exigem que os prefixos reais refinem o protocolo de referência e que seu sucesso esteja completo antes da publicação. Provas de ciclo de vida cobrem a correspondência entre voucher e recibo, conflitos, cancelamento e retirada. O teste de resposta perdida mata o serviço após persistir, reabre sem semente e confirma que a repetição não faz novas escritas. Fonte: applications/vouchers/README.md.}
\end{frame}

% 13
\begin{frame}{E1--E2: admissão e reutilização}
\begin{columns}[T,onlytextwidth]
\column{.52\textwidth}
\textbf{Fases ordinárias (s)}\par\vspace{2mm}
\begin{tikzpicture}
\begin{axis}[width=\linewidth,height=4cm,ybar,bar width=8pt,
 symbolic x coords={Leitura,Índice,Voucher},xtick=data,
 ymin=0,enlarge x limits=.23,grid=major,
 legend style={at={(.5,1.02)},anchor=south,legend columns=2,draw=none}]
\addplot[fill=accent,draw=accent] coordinates {
 (Leitura,\OrdinaryReadCheckingSeconds)
 (Índice,\OrdinaryIndexedCheckingSeconds)
 (Voucher,\OrdinaryVoucherCheckingSeconds)};
\addlegendentry{Verificação}
\addplot[fill=ink,draw=ink] coordinates {
 (Leitura,\OrdinaryReadCompilationSeconds)
 (Índice,\OrdinaryIndexedCompilationSeconds)
 (Voucher,\OrdinaryVoucherCompilationSeconds)};
\addlegendentry{Geração/compilação C}
\end{axis}
\end{tikzpicture}
\column{.45\textwidth}
\textbf{Chamada preparada (ms)}\par\vspace{4mm}
\small
\begin{tabular}{@{}lrr@{}}\toprule
Mediana & Volátil & Durável\\\midrule
Leitura & \numpt{\PreparedReadVolatileMedianMs} & \numpt{\PreparedReadPersistentMedianMs}\\[2mm]
Índice & \numpt{\PreparedIndexedVolatileMedianMs} & \numpt{\PreparedIndexedPersistentMedianMs}\\[2mm]
Voucher & \numpt{\PreparedVoucherVolatileMedianMs} & \numpt{\PreparedVoucherPersistentMedianMs}\\\bottomrule
\end{tabular}
\par\vspace{4mm}
Registro volátil da leitura:\\
\numpt{\RegistrationReadVolatileMedianSeconds} s, uma vez.
\end{columns}
\vspace{1mm}
\small
Worker: \numpt{\OrdinaryReadWorkerSharePercent}\% do tempo ordinário da leitura.\par
Registro amortizado em \BreakEvenReadCalls{} chamadas voláteis no modelo de medianas.
\source{E1: $n=\OrdinarySamples$. E2: $n=\PreparedSamples$ por programa/modo, \RegistrationSamples{} registros.}
\note{As barras mostram medianas de dois grupos do caminho ordinário: os dois checks do Book inteiro e a geração mais compilação C. Elas não somam toda a latência, pois há auditorias, staging e comunicação. As medianas totais são \numpt{\OrdinaryReadMedianSeconds}, \numpt{\OrdinaryIndexedMedianSeconds} e \numpt{\OrdinaryVoucherMedianSeconds} segundos. A lei generated\_oracle\_sound está dentro dos checks e sua parcela individual não pôde ser separada. O worker representa apenas \numpt{\OrdinaryReadWorkerSharePercent} por cento na leitura. A tabela à direita mostra chamadas preparadas em milissegundos, sem repetir verificação ou compilação. O registro paga esses custos uma vez. A fórmula é o teto de R dividido por O menos P, com unidades iguais. Os equilíbrios voláteis são \BreakEvenReadCalls, \BreakEvenIndexedCalls{} e \BreakEvenVoucherCalls{} chamadas para leitura, índice e voucher. São estimativas com medianas, excluindo inicialização. As operações têm um aquecimento, mas os registros não. O índice alterna o valor e o voucher usa IDs novos com retirada fora do intervalo da emissão. As distribuições do índice e voucher são amplas. O artigo separa throughput ativo em blocos do throughput completo, que inclui validações e retiradas.}
\end{frame}

% 14
\begin{frame}{E3: leitura durável tem a maior diferença}
\centering
\begin{tikzpicture}
\begin{axis}[width=.94\linewidth,height=4.5cm,ybar,bar width=10pt,
 ymode=log,log origin=infty,ymin=.01,ymax=120,
 ylabel={Mediana (ms), escala log},
 ytick={.01,.1,1,10,100},yticklabel={\pgfmathparse{exp(\tick)}\pgfmathprintnumber[fixed,precision=2,use comma]{\pgfmathresult}},
 symbolic x coords={L,U,PL,PU},xtick=data,
 xticklabels={{Pontual: leitura},{Pontual: atualização},{Proprietário: leitura},{Proprietário: atualização}},
 enlarge x limits=.15,grid=major,
 legend style={at={(.5,1.02)},anchor=south,legend columns=3,draw=none},
 nodes near coords, point meta=explicit symbolic, nodes near coords style={font=\scriptsize,fill=white,inner sep=.4pt}]
\addplot[fill=accent,draw=accent] coordinates {
 (L,\JayPointReadMedianMs)[\latpt{\JayPointReadMedianMs}] (U,\JayPointUpdateMedianMs)[\latpt{\JayPointUpdateMedianMs}]
 (PL,\JayOwnerReadMedianMs)[\latpt{\JayOwnerReadMedianMs}] (PU,\JayOwnerUpdateMedianMs)[\latpt{\JayOwnerUpdateMedianMs}]};
\addlegendentry{JayDB2 preparado}
\addplot[fill=ink,draw=ink] coordinates {
 (L,\PostgresPointReadMedianMs)[\latpt{\PostgresPointReadMedianMs}] (U,\PostgresPointUpdateMedianMs)[\latpt{\PostgresPointUpdateMedianMs}]
 (PL,\PostgresOwnerReadMedianMs)[\latpt{\PostgresOwnerReadMedianMs}] (PU,\PostgresOwnerUpdateMedianMs)[\latpt{\PostgresOwnerUpdateMedianMs}]};
\addlegendentry{PostgreSQL 16}
\addplot[fill=muted!35,draw=muted] coordinates {
 (L,\SqlitePointReadMedianMs)[\latpt{\SqlitePointReadMedianMs}] (U,\SqlitePointUpdateMedianMs)[\latpt{\SqlitePointUpdateMedianMs}]
 (PL,\SqliteOwnerReadMedianMs)[\latpt{\SqliteOwnerReadMedianMs}] (PU,\SqliteOwnerUpdateMedianMs)[\latpt{\SqliteOwnerUpdateMedianMs}]};
\addlegendentry{SQLite sem restrição}
\end{axis}
\end{tikzpicture}
\par\vspace{1mm}\small\raggedright
Leitura pontual: \ratiopt{\JayPointReadOverPostgres}$\times$ PostgreSQL.
Atualização: \ratiopt{\JayPointUpdateOverPostgres}$\times$.\par
Lei do administrador e FORCE RLS recusaram acesso entre usuários.
\source{Durável, $n=\ComparisonSamples$ por operação. Processo e estado completos a cada chamada JayDB2.}
\note{O eixo vertical é logarítmico e começa acima de zero. As barras são medianas em milissegundos, com um aquecimento e \ComparisonSamples{} chamadas por condição. Leitura e atualização são medidas tanto sem política de propriedade quanto com proprietário. SQLite continua sem essa restrição em ambos. As leituras JayDB2 são ordens de grandeza mais lentas que as SQL. As atualizações têm diferença menor: no perfil pontual, \ratiopt{\JayPointUpdateOverPostgres} vezes PostgreSQL e \ratiopt{\JayPointUpdateOverSqlite} vezes SQLite. Não podemos generalizar ordens de grandeza para todas as operações. PostgreSQL usa fsync, synchronous\_commit e full\_page\_writes ligados. SQLite usa WAL e synchronous FULL. JayDB2 lança um worker nativo novo, transfere o estado completo e sincroniza uma imagem completa mesmo nas leituras. Chamadas preparadas não repetem prova, então verificação não explica esta diferença. Não isolamos a contribuição de cada mecanismo. O intervalo JayDB2 termina após ler a resposta completa, antes de decodificá-la. SQL inclui o driver em execute e fetchall. JayDB2 atualiza com confirmação e leitura posterior fora do intervalo, enquanto SQL retorna o valor na atualização. Na política comum, JayDB2 usa tabelas do proprietário e uma lei administrativa. PostgreSQL usa coluna de dono e FORCE ROW LEVEL SECURITY, com papel sem superusuário ou BYPASSRLS. Ambos impediram leitura e alteração entre usuários, por rejeição de admissão ou por nenhuma linha afetada. A fonte dos resultados é comparison.json. A documentação da política SQL está em PostgreSQL 16, ddl-rowsecurity.}
\end{frame}

% 15
\begin{frame}{E4: verificação e tamanho das provas}
\begin{columns}[T,onlytextwidth]
\column{.51\textwidth}
\textbf{Tempo de parede (s)}\par\vspace{2mm}
\begin{tikzpicture}
\begin{axis}[width=\linewidth,height=4.15cm,ybar,bar width=24pt,
 symbolic x coords={Contrato,Completa},xtick=data,ymin=0,
 enlarge x limits=.5,grid=major,
 nodes near coords, nodes near coords style={font=\small,/pgf/number format/.cd,fixed,precision=1,use comma}]
\addplot[fill=accent,draw=accent] coordinates {
 (Contrato,\VerifyContractSeconds) (Completa,\VerifySeconds)};
\end{axis}
\end{tikzpicture}
\column{.46\textwidth}
\textbf{Linhas não vazias}\par\vspace{4mm}
\small
\begin{tabular}{@{}lr@{}}\toprule
Implementação \texttt{src/} & \numpt{\SrcNonblankLines}\\[2mm]
Especificação \texttt{contract/} & \numpt{\ContractNonblankLines}\\[2mm]
Provas \texttt{proof/} & \numpt{\ProofNonblankLines}\\[2mm]
Host \texttt{server/} & \numpt{\ServerNonblankLines}\\\bottomrule
\end{tabular}
\par\vspace{4mm}
Provas: \ratiopt{\ProofToImplementationRatio}$\times$ as linhas\\
da implementação Bend.
\end{columns}
\small Ambos os alvos passaram. Verificação completa: \NegativeFixtures{} controles negativos.
\source{Uma execução por alvo. Linhas incluem comentários e excluem documentação e código gerado.}
\note{As barras medem uma execução local de make verify-contract e uma de make verify, em sequência e sem limpeza deliberada de cache. São \numpt{\VerifyContractSeconds} e \numpt{\VerifySeconds} segundos, não latências de requisição. Os dois modos oficiais verificam testemunhas positivas. O alvo completo exige imports mantidos e \NegativeFixtures{} controles negativos com falha exatamente na definição false. make test também passou e levou \numpt{\TestSeconds} segundos. A tabela conta linhas físicas não vazias, incluindo comentários, nas áreas nomeadas. O quociente entre proof e src é \ratiopt{\ProofToImplementationRatio}. Ele mostra o tamanho substancial do artefato de prova, sem medir cobertura, força lógica, tempo de autoria ou manutenibilidade. As aplicações não fazem parte desse inventário. Há \numpt{\ProofLawDeclarations} leis, \numpt{\ProofDefinitions} definições e \ProofTheoremWitnesses{} testemunhas de topo com nome terminado em theorem. São contagens sintáticas distintas.}
\end{frame}

% 16
\begin{frame}{E5: escala do estado completo}
\centering
\begin{tikzpicture}
\begin{axis}[width=.91\linewidth,height=4.45cm,
 xmin=0,ymin=0,grid=major,
 xlabel={Estado canônico (kB decimais)},ylabel={Mediana (ms)},
 legend style={at={(.5,1.02)},anchor=south,legend columns=2,draw=none},
 axis lines=left]
\addplot[accent,thick,mark=*] coordinates {\ScalingVolatileReadCoordinates};
\addlegendentry{Leitura volátil}
\addplot[ink,thick,mark=square*] coordinates {\ScalingVolatileUpdateCoordinates};
\addlegendentry{Atualização volátil}
\addplot[accent,thick,dashed,mark=triangle*] coordinates {\ScalingPersistentReadCoordinates};
\addlegendentry{Leitura durável}
\addplot[ink,thick,dashed,mark=diamond*] coordinates {\ScalingPersistentUpdateCoordinates};
\addlegendentry{Atualização durável}
\end{axis}
\end{tikzpicture}
\par\vspace{1mm}\small\raggedright
Teto de \numpt{\StateLimitBytes} bytes: \numpt{\StateMaximumRows} linhas.\par
\numpt{\StateOversizedRows} linhas exigiriam \numpt{\StateOversizedKilobytes} kB, fora desta configuração.
\source{$n=\ScalingSamples$ por operação/tamanho/modo. Distribuições amplas e durável não monotônico.}
\note{As linhas conectam medianas observadas em \ScalingSmallRows, \ScalingMiddleRows, \numpt{\ScalingLargeRows} e \numpt{\StateMaximumRows} linhas. O eixo horizontal usa o tamanho canônico real em kB decimais. As linhas contínuas são voláteis e as tracejadas são duráveis, com leitura e atualização separadas. O artigo apresenta os quartis, que são amplos. Há um aumento marcado entre \ScalingSmallRows{} e \ScalingMiddleRows{} linhas. Apenas descrevemos esse salto, sem atribuir causa. A série durável de leitura cai do penúltimo ponto ao último, portanto não é monotônica. O benchmark mede também lookup linear na última chave, decodificação, validação e publicação, além da transferência. O teto de \numpt{\StateLimitBytes} bytes comporta \numpt{\StateMaximumRows} linhas com \numpt{\StateMaximumBytes} bytes. A linha seguinte retorna limit nos dois modos e preserva a raiz. A população maior exige \numpt{\StateOversizedBytes} bytes e não recebeu latência. A carga foi construída por uma transação certificada a partir de uma semente vazia. Tentativas iniciais de semente grande esgotaram combustível em Seed.initial\_valid e não são amostras de operação. Os maiores grupos foram medidos novamente após corrigir a sonda de limite, e os menores foram retidos. Não extraímos uma lei geral de escalabilidade dessas séries.}
\end{frame}

% 17
\begin{frame}{Limitações e ameaças à validade}
\begin{itemize}
\item Provas manuais e aplicações pequenas, construídas pelo autor.
\item Transferência de estado completo e processos novos por chamada.
\item Um cliente serial e uma configuração local de medição.
\item Caminhos de armazenamento e respostas distintos entre bancos.
\end{itemize}
\gap
{\color{accent}A garantia aplica as leis instaladas sob T1--T3.}
\note{O administrador ainda precisa especificar a política desejada. Uma lei ausente não vira uma obrigação automaticamente. O exemplo do índice tem dois usuários fixos e uma linha booleana por usuário. Os recibos protegem o ID enquanto forem retidos e o estado autoritativo sobreviver. As medições usam uma máquina local, sem limpeza deliberada de cache ou repetição independente. Cache, armazenamento, sincronização e posicionamento de CPU podem mudar os resultados. Os hooks de fase afetam E1 e E2. Chamadas na mesma rodada são correlacionadas. Ordinário e preparado usam níveis distintos de otimização C, então sua diferença mede os perfis de serviço. Uma leitura bem-sucedida também pode seguir a publicação durável de uma raiz inalterada. Canais de tempo e consumo físico ficam fora do modelo de observações.}
\end{frame}

% 18
\begin{frame}{Contribuições e próximos passos}
\textbf{Resultado}\par\vspace{3mm}
Certificados ligados à execução e à publicação sob T1--T3.\par\vspace{3mm}
Leitura: \numpt{\OrdinaryReadMedianSeconds} s ordinária,
\numpt{\PreparedReadVolatileMedianMs} ms preparada volátil.\par\vspace{3mm}
Reuso amortizado em \BreakEvenReadCalls{} chamadas no modelo de medianas.
\gap
\textbf{Próximos passos}\par\vspace{3mm}
Provar durabilidade após falhas, reduzir T3 e o custo de estado completo.
\source{Fundamentos: Necula e Lee (1996), Ynot (2008), Malecha et al. (2010), Fiat (2015).}
\note{A contribuição principal conecta a prova do cliente às funções que iniciam a máquina, acompanham os prefixos e selecionam o resultado publicado. O contrato é habitado e exclui comportamentos errados concretos. O controlador preparado participa do fold verificado e as aplicações mostram permissões e invariantes distintos. A avaliação quantifica o custo: segundos na admissão ordinária e milissegundos no reuso preparado, com equilíbrio volátil estimado em \BreakEvenReadCalls{} chamadas para leitura. Mesmo preparada, a leitura pontual durável custa \ratiopt{\JayPointReadOverPostgres} vezes PostgreSQL. As provas têm \ratiopt{\ProofToImplementationRatio} vezes as linhas da implementação Bend e o teto atual admite \numpt{\StateMaximumRows} linhas da carga de escala. Os próximos passos incluem prova de durabilidade e recuperação, menos confiança física em T3, admissão mais rápida e menor transferência e publicação de estado completo. As referências completas estão no artigo.}
\end{frame}

\appendix
\begin{frame}{Apêndice: teoremas de execução}
\small
\begin{tabular}{@{}L{.31\linewidth}L{.62\linewidth}@{}}
\textbf{StartupTheorem} & Inicialização executada e verificação de recuperação.\\[3mm]
\textbf{InitTheorem} & Auxiliar abstrato de admissão da semente para Fold.\\[3mm]
\textbf{CoreTheorem} & Admitted, Safe, Completed e Progressing.\\[3mm]
\textbf{IngressTheorem} & Decodificação ordinária e raiz exata.\\[3mm]
\textbf{UnifiedCoreTheorem} & Seleção direta/registrada do pedido exato.\\[3mm]
\textbf{FoldTheorem} & Validade em toda sequência finita de transações.\\
\end{tabular}
\source{Testemunhas: Configuration, Init, CoreClauses, Ingress e Iteration.}
\note{Startup fala de ConfigurationCheck.initialize e recover, usados pelo serviço. A declaração exata de validade da semente já deve ter sido admitida sob T1. Init é o auxiliar abstrato que fornece a premissa do fold, e não o inicializador executado. Core reúne as quatro cláusulas com as mesmas premissas. Ingress cobre o corpo ordinário completo e UnifiedCore a seleção direta ou preparada real. Fold deriva a validade de cada raiz alcançada a partir da semente aceita e de transições computadas. Essas diferenças evitam confundir um auxiliar lógico com a implementação de startup.}
\end{frame}

\begin{frame}{Apêndice: codificação e controles}
\small
\begin{tabular}{@{}L{.35\linewidth}L{.58\linewidth}@{}}
\textbf{CodecRoundtrip} & Codificar e decodificar preserva o valor nos limites.\\[2mm]
\textbf{CodecCanonical} & Bytes aceitos recodificam exatamente.\\[2mm]
\textbf{FramingRoundtrip} & Inversão do pacote privado de dois blocos.\\[2mm]
\textbf{ControlTheorem} & Autenticação, recibos limitados e deltas exatos.\\[2mm]
\textbf{ServiceFoldTheorem} & Raízes válidas e histórico exato do registro.\\[2mm]
\textbf{ServiceNoRevivalTheorem} & Handle alocado e ausente permanece ausente.\\
\end{tabular}
\source{Pontes: WorkerBridge, CertifiedWorker, WorkerDecision, DispatchBridge e ControlRegistryBytes.}
\note{O roundtrip de valores exige que o tamanho seja representável e caiba no teto do decodificador. FramingRoundtrip trata do pacote privado, enquanto o quadro público é Protocol.frame. Control prova a estrutura das decisões mesmo sem assumir a correção de um token de admissão arbitrário. ServiceFold acrescenta essa premissa T1 para derivar certificação universal das fontes vivas. ServiceNoRevival mantém a ausência do handle ao longo do restante da execução. As pontes relacionam o worker e a decisão emitida à composição, e os bytes do registro do controlador ao decodificador preparado. Fontes: proof/README.md e contract/Theorems.bend.}
\end{frame}

\begin{frame}{Apêndice: teoremas do registro}
\small
\begin{tabular}{@{}L{.40\linewidth}L{.53\linewidth}@{}}
\textbf{RegistryLookupTheorem} & Lookup vivo seleciona uma fonte universalmente certificada.\\[5mm]
\textbf{RegistryInvalidationTheorem} & Remoção elimina lookup e conserva histórico de alocação.\\[5mm]
\textbf{ServiceLifetimeTheorem} & Handle já alocado continua reservado e nova preparação recusa sua reutilização.\\
\end{tabular}
\source{Componentes públicos de ServiceFoldTheorem. Testemunhas: ControlHistory.}
\note{ServiceFold inclui estes três resultados além da validade das raízes e do histórico exato do registro. RegistryLookup extrai certificação universal da fonte encontrada em um registro acompanhado pelas provas de admissão. RegistryInvalidation elimina o lookup no escopo da instância e preserva os handles usados. ServiceLifetime mantém o handle alocado em todo ponto futuro e mostra que preparar esse candidato recusa. ServiceNoRevival combina essa reserva com a ausência do lookup, preservando a revogação em qualquer sequência posterior de transações e controles. A execução física recebe os bytes atuais sob T3.}
\end{frame}

\begin{frame}{Apêndice: T3, entrada e execução}
\begin{tabular}{@{}L{.08\linewidth}L{.85\linewidth}@{}}
\textbf{1} & Congelar configuração, fontes, ferramentas e raiz inicial válida.\\[4mm]
\textbf{2} & Ler quadros limitados, preservar EOF e fornecer o corpo e a raiz atuais.\\[4mm]
\textbf{3} & Compilar e carregar o Book exato da tupla admitida.\\[4mm]
\textbf{4} & Executar worker privado novo pelo caminho verificado, sem poder de publicação.\\
\end{tabular}
\source{Numeração normativa: proof/README.md, T3.1--T3.4.}
\note{Esses são os quatro primeiros deveres físicos do host. A configuração e os bytes das ferramentas ficam congelados. O corpo original completo e um único snapshot atual atravessam a fronteira de processos, com os bytes extras da extensão contados no limite. A tupla admitida mantém código, prova, contexto e argumentos até o lançamento do artefato correspondente. O worker é novo e só manipula a proposta privada. A ordem lógica de observações deve permanecer válida mesmo quando há interrupção. A prova Bend especifica os valores e as transições, e estes deveres ligam a entrada física àqueles valores.}
\end{frame}

\begin{frame}{Apêndice: T3, publicação e registro}
\begin{tabular}{@{}L{.08\linewidth}L{.85\linewidth}@{}}
\textbf{5} & Recolher saída completa e concluir limpeza antes de publicar.\\[3mm]
\textbf{6} & Publicar atomicamente a raiz e resposta de Commit. Keep preserva a raiz.\\[3mm]
\textbf{7} & Encaminhar os últimos bytes do registro sem alteração, a tupla T1 exata e instalar o delta verificado.\\[3mm]
\textbf{8} & Entregar os bytes exatos. Falha após commit mantém a publicação.\\
\end{tabular}
\source{O registro canônico não é autenticado. Revogação depende do encaminhamento atual e exato.}
\note{Antes do commit, o host termina a coleta, espera o processo e faz a limpeza que ainda permite manter a raiz anterior. Depois obedece exatamente à soma Commit ou Keep. No controle preparado, o host fornece o registro mais recente e a tupla universal recebida da admissão, conserva o histórico de handles e gerencia o artefato correspondente. Recibos de controle preservam a raiz e a última resposta publicada, e uma falha de entrega não desfaz um controle. Bytes antigos ou fabricados podem quebrar revogação e a associação do handle. A certificação universal do programa executado continua dependendo de T1 e do matching gerado.}
\end{frame}

\begin{frame}{Apêndice: T3, erros e respostas Rust}
\begin{tabular}{@{}L{.08\linewidth}L{.85\linewidth}@{}}
\textbf{9} & Tabela de códigos: erros de admissão/prova dão \texttt{certificate}; demais códigos dão \texttt{failure}.\\[4mm]
\textbf{10} & Artefato preparado ausente: manter raiz e retornar o quadro fixo de falha.\\[4mm]
\textbf{11} & Integridade do artefato inválida: mesma resposta, antes de executar.\\
\end{tabular}
\gap
\small As duas respostas exclusivas de Rust são comparadas byte a byte com a negação verificada de handle ausente.
\note{Uma prova inválida rejeita a admissão do Book inteiro antes de Gate executar. A classificação física usa a tabela exata de proof/README.md, T3.9: IMPORT, IMPORT\_CYCLE, DEPENDENCY, PARSE, META, UNSAFE, FOREIGN, IO, OPEN\_LAW, SHADOW, REDEFINITION, PRIMARY\_CHECK, BENDTT\_CHECK, TARGET, SOURCE\_BINDING, HOST\_CAPTURE e KERNEL\_CAPTURE viram certificate. Outros códigos, incluindo HOST\_COMPILE após admissão pura bem-sucedida, viram failure. HANDLE e PREPARED\_INTEGRITY geram as duas respostas diretamente em Rust. Os testes verificam sua igualdade byte a byte com a negação do worker. Artefato indisponível para registro vivo é falha de disponibilidade.}
\end{frame}
\end{document}
